\documentclass[11pt,a4paper]{article}

\usepackage[utf8]{inputenc}
\usepackage[T1]{fontenc}
\usepackage{amsmath,amssymb,amsfonts}
\usepackage{physics}
\usepackage{graphicx}
\usepackage{booktabs}
\usepackage{array}
\usepackage{siunitx}
\usepackage{hyperref}
\usepackage[margin=1in]{geometry}
\usepackage{xcolor}
\usepackage{listings}
\usepackage{enumitem}
\usepackage{fancyhdr}
\usepackage{titlesec}

\hypersetup{
    colorlinks=true,
    linkcolor=blue!60!black,
    urlcolor=blue!60!black,
    citecolor=blue!60!black,
    pdftitle={PY-BEMCS Physics Reference and User Manual},
    pdfauthor={Dr. Bharat Singh Rawat}
}

\pagestyle{fancy}
\fancyhf{}
\fancyhead[L]{PY-BEMCS Reference Manual}
\fancyhead[R]{\thepage}
\fancyfoot[C]{\small Python Beam Extraction \& Monte Carlo Simulator}

\titleformat{\section}{\Large\bfseries\color{blue!50!black}}{\thesection}{1em}{}
\titleformat{\subsection}{\large\bfseries\color{blue!40!black}}{\thesubsection}{1em}{}
\titleformat{\subsubsection}{\normalsize\bfseries\color{blue!30!black}}{\thesubsubsection}{1em}{}

\newcommand{\PYBEMCS}{PY-BEMCS}

\lstset{
    basicstyle=\small\ttfamily,
    backgroundcolor=\color{gray!10},
    frame=single,
    breaklines=true,
    captionpos=b
}

\begin{document}

% ---------------------- TITLE ----------------------
\begin{titlepage}
\centering
\vspace*{2cm}

{\Huge\bfseries Python Beam Extraction \&\\[0.3em] Monte Carlo Simulator \par}
\vspace{0.5cm}
{\Large\bfseries (\PYBEMCS)\par}

\vspace{1.2cm}
{\large\textit{Comprehensive Physics Reference and User Manual}\par}

\vspace{2.5cm}
{\large Dr.~Bharat Singh Rawat\par}
\vspace{0.2cm}
{\normalsize \href{mailto:bharat.bharat22@gmail.com}{bharat.bharat22@gmail.com}\par}

\vspace{1.5cm}
{\normalsize \textbf{Version 2.0 (Extended Modular Release)}\par}
\vspace{0.2cm}
{\normalsize \today\par}

\vfill
\hrule
\vspace{0.5em}
{\small Open-Source 2D3V Electrostatic Particle-in-Cell, Plume MCC, Thermo-Mechanical Grid Erosion, \& Multi-Core Benchmarking Suite}
\end{titlepage}

% ---------------------- ABSTRACT ----------------------
\begin{center}
{\Large\bfseries Abstract}
\end{center}
\vspace{0.5em}

\PYBEMCS{} is an open-source 2D3V (two spatial, three velocity components) electrostatic Particle-in-Cell (PIC) and Monte Carlo Collision (MCC) simulation suite engineered for ion thruster optics, neutral beam injectors (NBI), and advanced plasma extraction systems. This document provides a complete mathematical, physical, and numerical specification of the simulator, together with comprehensive user guidelines for the graphical interface, headless batch solvers, and automated multi-core benchmarking scripts.

Key capabilities detailed in this manual include:
\begin{itemize}[nosep]
    \item Self-consistent Bohm extraction from upstream plasma sheaths with space-charge Poisson solves and Boltzmann electron screening.
    \item 2D3V relativistic Boris particle advancement with dual Taichi GPU/CPU accelerated and pure-NumPy fallback kernels.
    \item Charge-Exchange (CEX) collision dynamics using Roy's neutral plume model or user-imported experimental cross-section datasets with log-log cubic spline smoothing.
    \item In-situ sputter erosion modeling (Matsunami yield formulation) coupled with live grid mesh boundary deletion.
    \item Coupled 2D thermal conduction, Stefan--Boltzmann radiative cooling, and Euler--Bernoulli cantilever thermo-mechanical grid aperture deflection.
    \item Modular JSON configuration architecture (\texttt{config.json}) supporting arbitrary $N$-grid configurations, discharge chamber plasma dynamics, and live parameter reloading.
    \item Automated multi-core parallel benchmarking suite for CEX erosion rates, Electron Backstreaming (EBS) limits, primary beam impingement / grid scraping, and perveance-divergence curves.
    \item Built-in asynchronous standalone application packager (\texttt{PyInstallerWorker}) embedded directly in the GUI.
\end{itemize}

\vspace{1em}
\hrule
\vspace{1em}

\tableofcontents
\newpage

% =====================================================
\section{Computational Framework and Architecture}
% =====================================================

\subsection{Coordinate System and Discretisation}

The computational domain employs a 2D Cartesian grid in $(x, y)$, where $x$ represents the primary axial beam extraction coordinate and $y$ corresponds to the transverse/radial coordinate (half-aperture or axisymmetric slice). The domain extents are defined by $L_x \times L_y$ with uniform cell spacing $\Delta x \times \Delta y$:
\begin{equation}
    n_x = \left\lfloor\frac{L_x}{\Delta x}\right\rfloor + 1, \qquad
    n_y = \left\lfloor\frac{L_y}{\Delta y}\right\rfloor + 1.
\end{equation}
By default, $\Delta x = \Delta y = \SI{0.02}{\milli\meter}$ ($\SI{20}{\micro\meter}$), ensuring the grid spacing resolves the local Debye length throughout the sheath region while preserving computational tractability.

Every macroparticle stores three-dimensional velocity components $\vb{v} = (v_x, v_y, v_z)$. While spatial motion is constrained to the $(x, y)$ plane, $v_z$ is rigorously advanced in time during magnetic rotations and collisional scattering.

\subsection{Dual-Kernel Execution Engine (Taichi GPU + CPU Fallback)}

To combine peak GPU performance on workstation hardware with maximum portability across containerized or headless computing clusters, \PYBEMCS{} implements a dual-kernel architecture:
\begin{enumerate}[nosep]
    \item \textbf{Taichi Backend (GPU/CPU)}: When Taichi is initialized, particle pushes, bilinear charge deposition, and explicit thermal conduction compile into parallel GPU kernels (CUDA, Metal, or Vulkan) running in 32-bit floating-point precision ($\texttt{ti.f32}$).
    \item \textbf{Vectorised CPU Fallback Kernels}: If Taichi is not present (or when running inside frozen executable packages), the physics engine automatically activates vectorised CPU routines (\texttt{accumulate\_rho\_cpu}, \texttt{push\_particles\_boris\_cpu}) using optimized NumPy array operations.
\end{enumerate}

% =====================================================
\section{Particle-in-Cell Electrostatic Algorithm}
% =====================================================

\subsection{Self-Consistent Poisson--Boltzmann Solver}

The electrostatic potential $\phi(\vb{x})$ is governed by Poisson's equation including macroparticle ion density and thermalized Boltzmann electrons:
\begin{equation}
    \nabla^{2}\phi(\vb{x}) \;=\; -\,\frac{1}{\varepsilon_{0}}
        \left[ q_{\mathrm{i}}\,n_{\mathrm{i}}(\vb{x}) - e\,n_{\mathrm{e}}(\vb{x}) \right],
    \label{eq:poisson}
\end{equation}
where $q_{\mathrm{i}} = Z e$ is the ion charge, $n_{\mathrm{i}}$ is the deposited ion macroparticle number density, and the electron density $n_{\mathrm{e}}$ follows the non-linear Boltzmann relation:
\begin{equation}
    n_{\mathrm{e}}(\vb{x}) \;=\; n_{0}\,
    \exp\!\left[\frac{\min(\phi(\vb{x}),\,\phi_{\mathrm{pl}}) - \phi_{\mathrm{pl}}}{T_{\mathrm{e}}}\right],
    \label{eq:boltzmann}
\end{equation}
with $T_{\mathrm{e}}$ in electron-volts ($\si{\eV}$), $n_0$ the upstream plasma density, and $\phi_{\mathrm{pl}}$ the reference plasma potential established at the upstream boundary ($x = 0$). The clipping operator $\min(\phi, \phi_{\mathrm{pl}})$ prevents non-physical exponential density blow-ups in regions where local potential excursions exceed the upstream plasma potential.

Equation~\eqref{eq:poisson} is discretised on the uniform mesh via a standard 5-point finite-difference stencil:
\begin{equation}
    \frac{\phi_{i+1,j} - 2\phi_{i,j} + \phi_{i-1,j}}{(\Delta x)^2} +
    \frac{\phi_{i,j+1} - 2\phi_{i,j} + \phi_{i,j-1}}{(\Delta y)^2}
    = -\frac{\rho_{i,j}}{\varepsilon_{0}}.
\end{equation}
The resulting sparse linear system $A \boldsymbol{\phi} = \mathbf{b}$ is factorised using SuperLU sparse LU decomposition. Non-linear electron space charge is converged via successive under-relaxation with relaxation parameter $\omega = 0.2$.

\subsection{Bilinear Charge Deposition}

Macroparticle charge is mapped to the four surrounding mesh nodes using bilinear interpolation weights. For a particle located at $(x_p, y_p)$ with grid indices $i = \lfloor x_p/\Delta x \rfloor$, $j = \lfloor y_p/\Delta y \rfloor$:
\begin{align}
    h_x &= \frac{x_p - i\Delta x}{\Delta x}, \qquad
    h_y = \frac{y_p - j\Delta y}{\Delta y},\\
    \rho_{i,j}     &\mathrel{+}= \frac{q_{\mathrm{i}} w}{V_{\mathrm{cell}}}\,(1-h_x)(1-h_y),\\
    \rho_{i+1,j}   &\mathrel{+}= \frac{q_{\mathrm{i}} w}{V_{\mathrm{cell}}}\,h_x(1-h_y),\\
    \rho_{i,j+1}   &\mathrel{+}= \frac{q_{\mathrm{i}} w}{V_{\mathrm{cell}}}\,(1-h_x)h_y,\\
    \rho_{i+1,j+1} &\mathrel{+}= \frac{q_{\mathrm{i}} w}{V_{\mathrm{cell}}}\,h_x h_y,
\end{align}
where $V_{\mathrm{cell}} = \Delta x \Delta y \Delta z$ is the cell volume with assumed unit transverse thickness $\Delta z = \SI{1}{\milli\meter}$, and $w$ is the statistical macroparticle weighting factor:
\begin{equation}
    w = \max\!\left(\frac{n_0 V_{\mathrm{cell}}}{N_{\mathrm{ppc}}},\; 10^{3}\right),
\end{equation}
with default nominal target $N_{\mathrm{ppc}} = 40$ particles per cell.

\subsection{Electric Field Computation and Boris Pusher}

Electric field components are evaluated at mesh nodes via second-order central differencing:
\begin{equation}
    E_{x,i,j} = -\frac{\phi_{i+1,j} - \phi_{i-1,j}}{2\Delta x}, \qquad
    E_{y,i,j} = -\frac{\phi_{i,j+1} - \phi_{i,j-1}}{2\Delta y},
\end{equation}
and interpolated to particle positions. Particle trajectories are advanced using the explicit leap-frog Boris method~\cite{birdsall}:
\begin{align}
    \vb{v}^{-} &= \vb{v}^{n-1/2} + \frac{q\,\vb{E}}{2m}\,\Delta t,\\
    \vb{t} &= \frac{q\,\vb{B}}{2m}\,\Delta t, \qquad
    \vb{s} = \frac{2\,\vb{t}}{1 + |\vb{t}|^2},\\
    \vb{v}' &= \vb{v}^{-} + \vb{v}^{-} \times \vb{t},\\
    \vb{v}^{+} &= \vb{v}^{-} + \vb{v}' \times \vb{s},\\
    \vb{v}^{n+1/2} &= \vb{v}^{+} + \frac{q\,\vb{E}}{2m}\,\Delta t,\\
    \vb{x}^{n+1} &= \vb{x}^{n} + \vb{v}^{n+1/2}\,\Delta t.
\end{align}
This scheme guarantees exact conservation of phase-space volume and energy conservation in static magnetic fields.

% =====================================================
\section{Plasma Injection and Bohm Optics}
% =====================================================

\subsection{Bohm Sheath Extraction Criterion}

Ions enter the computational domain through the upstream plasma boundary ($x = 0$) with an initial axial drift satisfying the Bohm sheath velocity criterion:
\begin{equation}
    v_{\mathrm{B}} = \sqrt{\frac{Z e T_{\mathrm{e}}}{m_{\mathrm{ion}}}}, \qquad
    m_{\mathrm{ion}} = M_{\mathrm{amu}} \cdot u,
    \label{eq:bohm}
\end{equation}
where $M_{\mathrm{amu}}$ is the ion mass in atomic mass units and $u = \SI{1.66054e-27}{\kilo\gram}$.

The injected ion current across the open sheath area $A_{\mathrm{inj}}$ is:
\begin{equation}
    I_{\mathrm{ion}} = 0.61\, Z e n_{0} v_{\mathrm{B}} A_{\mathrm{inj}},
\end{equation}
where the prefactor $0.61 \approx \exp(-1/2)$ accounts for the quasi-neutral pre-sheath density drop. The number of macroparticles injected at each timestep $\Delta t$ is:
\begin{equation}
    N_{\mathrm{inj}} = \frac{I_{\mathrm{ion}} \Delta t}{Z e w} + \xi_{\mathrm{stochastic}},
\end{equation}
with stochastic remainder accumulation ensuring exact long-term current conservation.

\subsection{Thermal Velocity Sampling}

Thermal spread at the sheath boundary is sampled from a 3D Maxwellian distribution at ion temperature $T_{\mathrm{i}}$:
\begin{equation}
    v_x = v_{\mathrm{B}} + \zeta_x v_{\mathrm{th,i}}, \qquad
    v_y = \zeta_y v_{\mathrm{th,i}}, \qquad
    v_z = \zeta_z v_{\mathrm{th,i}},
\end{equation}
where $\zeta_{x,y,z} \sim \mathcal{N}(0, 1)$ are standard normal random variables, and $v_{\mathrm{th,i}} = \sqrt{Z e T_{\mathrm{i}} / m_{\mathrm{ion}}}$.

% =====================================================
\section{Multiphysics Grid Models}
% =====================================================

\subsection{Arbitrary $N$-Grid Geometry Specification}

The extraction optics consist of an arbitrary sequence of $N$ conducting grids defined by:
\begin{itemize}[nosep]
    \item DC potential $V_k$ (\si{\volt})
    \item Thickness $t_k$ (\si{\milli\meter})
    \item Inter-grid axial gap to subsequent grid $g_k$ (\si{\milli\meter})
    \item Aperture radius $r_k$ (\si{\milli\meter})
    \item Upstream chamfer angle $\theta_k$ (\si{\degree})
\end{itemize}

Solid boundary masks $\mathcal{M}_k(x,y)$ are dynamically synthesised and rasterised onto the finite-difference mesh.

\subsection{Cantilever Thermo-Mechanical Grid Deformation}

Under intense ion impingement and radiative heating, individual grid membranes experience substantial temperature rises $\Delta T_k = T_k - T_{\infty}$. Each grid is modeled as an annular plate segment undergoing Euler--Bernoulli cantilever bending clamped at the outer boundary ($y = L_y$) and free at the aperture edge ($y = r_k$):
\begin{equation}
    \delta_{\max,k} = \frac{\alpha_k \Delta T_k L_k^2}{2 t_k}, \qquad L_k = L_y - r_k,
\end{equation}
where $\alpha_k$ is the linear thermal expansion coefficient ($\si{\per\kelvin}$). The transverse displacement profile:
\begin{equation}
    \delta(y) = \delta_{\max,k} \left(\frac{L_y - y}{L_k}\right)^2,
\end{equation}
is dynamically added as an axial coordinate distortion to the grid boundaries. Whenever cumulative structural deformation exceeds $\SI{5}{\micro\meter}$, the Poisson sparse matrix operator is automatically re-factored.

\subsection{2D Thermal Diffusion and Radiative Dissipation}

Thermal energy deposition from impacting ions is balanced by 2D conduction within the grid solid mask and Stefan--Boltzmann radiative cooling to space:
\begin{equation}
    \rho_k c_{p,k} \frac{\partial T}{\partial t} = k_k \left(\frac{\partial^2 T}{\partial x^2} + \frac{\partial^2 T}{\partial y^2}\right) + \dot{q}_{\mathrm{dep}} - \frac{\varepsilon_k \sigma_{\mathrm{SB}}}{\Delta z}\left(T^4 - T_{\infty}^4\right),
\end{equation}
where $k_k$ is thermal conductivity, $c_{p,k}$ is specific heat capacity, $\varepsilon_k$ is surface emissivity, $\sigma_{\mathrm{SB}} = \SI{5.67e-8}{\watt\per\meter\squared\per\kelvin\tothe{4}}$, and $T_{\infty} = \SI{300}{\kelvin}$.

% =====================================================
\section{Collisions and Sputter Erosion}
% =====================================================

\subsection{Charge-Exchange (CEX) Dynamics and Plume Neutral Model}

Neutral gas density distribution throughout the grid gap and downstream region follows Roy's analytical plume expansion model~\cite{roy}:
\begin{equation}
    n_{\mathrm{n}}(x,y) = n_{\mathrm{n},0} \, a \left[1 - \frac{1}{\sqrt{1 + (r_T / R)^2}}\right] \cos\theta,
\end{equation}
with $R = \sqrt{y^2 + (x + r_T)^2}$, $\theta = \arctan(y / (x + r_T))$, and normalisation factor $a = (1 - 1/\sqrt{2})^{-1}$.

The null-collision Monte Carlo collision (MCC) probability for an ion traversing distance $\Delta s = |\vb{v}|\Delta t$ is:
\begin{equation}
    P_{\mathrm{CEX}} = 1 - \exp\left[-n_{\mathrm{n}}(x,y) \, \sigma_{\mathrm{CEX}}(E) \, |\vb{v}| \, \Delta t\right].
\end{equation}
The analytical xenon CEX cross-section formulation is:
\begin{equation}
    \sigma_{\mathrm{CEX}}(g) = \left[-0.8821\ln(g) + 15.1262\right]^2 \times 10^{-20} \;\si{\meter\squared},
\end{equation}
where $g = |\vb{v}|$ is relative collision speed. If custom cross-section CSV tables are imported, this formula is dynamically superseded by cubic spline interpolation in $\log_{10}(E)$--$\log_{10}(\sigma)$ space.

\subsection{Sputter Yield and Dynamic Cell Removal}

Erosion of grid boundaries under ion bombardment is governed by the empirical Matsunami sputtering yield relation~\cite{matsunami}:
\begin{equation}
    Y(E) = \min\!\left[Y_0 \left(E_{\mathrm{eV}} - E_{\mathrm{th}}\right)^{1.5},\; 1.0\right],
\end{equation}
where $Y_0$ is the sputter yield coefficient and $E_{\mathrm{th}}$ is the sputtering threshold energy ($\si{\eV}$).

Cumulative damage per cell is accumulated as:
\begin{equation}
    D_{i,j} \mathrel{+}= Y(E) \cdot \kappa_{\mathrm{acc}},
\end{equation}
where $\kappa_{\mathrm{acc}}$ is the numerical time-acceleration factor. When $D_{i,j} \ge D_{\mathrm{fail}}$, the grid cell is removed from the solid mask, exposing new interior nodes and updating the field solver in real time.

% =====================================================
\section{Modular Configuration Architecture (\texttt{config.json})}
% =====================================================

The simulator parameters are managed through a unified, hierarchical JSON structure (\texttt{config.json}). A representative configuration is shown below:

\begin{lstlisting}[language=json,caption={Standard \texttt{config.json} structure}]
{
  "beam_species": {
    "mass_amu": 131.293,
    "charge_state": 1
  },
  "grid_material": {
    "preset": "Molybdenum"
  },
  "simulation": {
    "n0_plasma": 7.73e17,
    "Te_up": 8.5,
    "Ti": 0.034,
    "Tn": 400.0,
    "n0": 5.27e18,
    "Accel": 0.5,
    "Thresh": 10000.0,
    "sim_mode": "Both",
    "geometry": "half_hole"
  },
  "discharge_chamber": {
    "radius_m": 0.02,
    "length_m": 0.10,
    "n_cusp": 4,
    "pitch_mm": 3.0
  },
  "grids": [
    { "V": 1300.0, "t": 1.2, "gap": 0.9, "r": 1.3, "cham": 20.0 },
    { "V": -150.0, "t": 0.6, "gap": 0.5, "r": 0.8, "cham": 0.0 },
    { "V": 50.0,   "t": 0.5, "gap": 4.0, "r": 1.2, "cham": 0.0 }
  ],
  "cross_sections": {
    "cx_file": "",
    "see_file": "",
    "custom_file": "",
    "spline_smoothing": 0.0
  },
  "terminal_output": {
    "grid_temperatures": true,
    "beam_divergence": true,
    "saddle_point_potential": true,
    "mean_particle_energy": true,
    "iteration_time": true
  }
}
\end{lstlisting}

% =====================================================
\section{Graphical User Interface and Features}
% =====================================================

\subsection{Menu Bar and Advanced Controls}

\begin{itemize}
    \item \textbf{Settings $\rightarrow$ Advanced Parameters...}: Configures neutralizer injection coordinates $(x_{\mathrm{n}}, r_{\mathrm{n}})$, plasma offset potential, and electron mass scaling ratio $\mathcal{R} = m_{\mathrm{ion}}/m_{\mathrm{e}}$.
    \item \textbf{Settings $\rightarrow$ Open Config JSON...}: Interactive file browser to load external JSON configuration profiles.
    \item \textbf{Settings $\rightarrow$ Reload config.json}: Hot-reloads configuration files from disk without restarting the application.
    \item \textbf{Settings $\rightarrow$ Build .exe...}: Compiles the current Python source code into a standalone single-file binary executable using an asynchronous background thread with live milestone tracking.
    \item \textbf{Beam $\rightarrow$ Ion Species...}: Configures beam ion atomic mass (amu) and ionization charge state ($+1, +2, \dots$) with predefined presets (Xe, Kr, Ar, N$_2$, O$_2$, H$_2$, He, Hg, Cs).
    \item \textbf{Beam $\rightarrow$ Cross-Section Manager...}: Import external experimental CSV tables for CEX, SEE yield, or custom reaction channels with interactive log-log cubic spline fitting.
    \item \textbf{Materials $\rightarrow$ Grid Material...}: Selects material presets (Molybdenum, Steel SS316, Titanium, Graphite) or custom thermophysical properties ($k, \rho, c_p, \varepsilon, \alpha, E, Y_0, E_{\mathrm{th}}$).
\end{itemize}

\subsection{Live Diagnostics and Data Export}

\begin{itemize}[nosep]
    \item \textbf{Beam Divergence Angle ($\theta_{95}$)}: Evaluated from the 95\textsuperscript{th} percentile of ion trajectory angles downstream of the last grid:
    \begin{equation}
        \theta_{95} = Q_{0.95}\!\left[\,\left|\arctan\left(\frac{v_y}{v_x}\right)\right|\,\right].
    \end{equation}
    \item \textbf{Electron Backstreaming (EBS) Saddle Potential}: Continuous monitoring of centerline potential in the accelerator aperture to verify backstreaming suppression ($V_{\mathrm{sp}} < -5\,\si{\volt}$).
    \item \textbf{Ion Energy Distribution Function (IEDF)}: Live spectral histogram tracking particle energy distributions exiting the grid system.
    \item \textbf{Asynchronous Data Exporters}: Background progress dialogs for exporting time-series simulation telemetry to CSV, particle phase-space archives, and animated GIF recordings.
\end{itemize}

% =====================================================
\section{Multi-Core Benchmarking Suite (\texttt{benchmarks/})}
% =====================================================

The \texttt{benchmarks/} directory contains dedicated multi-core physics sweep scripts utilising Python's \texttt{multiprocessing} engine to execute parametric studies across independent CPU worker processes simultaneously:

\begin{table}[h!]
\centering
\small
\begin{tabular}{lp{6cm}p{4.5cm}}
\toprule
\textbf{Benchmark Script} & \textbf{Physics Objective} & \textbf{Generated Output} \\
\midrule
\texttt{benchmark\_cex.py} & Sputter erosion rate vs background neutral density ($n_{\mathrm{n}} = 10^{18}$--$5\times 10^{19}\,\si{\per\meter\cubed}$) across multiple grid gap configurations. & \texttt{cex\_benchmark\_multicore.png} \\
\addlinespace
\texttt{benchmark\_ebs.py} & Centerline saddle-point potential barrier vs accelerator grid voltage ($V_{\mathrm{a}} = -400$--$0\,\si{\volt}$) evaluating the EBS safety boundary. & \texttt{ebs\_benchmark\_multicore.png} \\
\addlinespace
\texttt{benchmark\_impingement.py} & Direct primary beam ion grid interception and scraping percentage vs perveance ($P = I_{\mathrm{ion}} / V_{\mathrm{tot}}^{1.5}$). & \texttt{impingement\_ratio\_benchmark.png} \\
\addlinespace
\texttt{benchmark\_perveance.py} & Beam divergence angle vs perveance over 60 plasma density increments across multiple grid gap spacings. & \texttt{divergence\_vs\_perveance\_benchmark\_multicore.png} \\
\addlinespace
\texttt{benchmark\_perveance\_Vs\_Sweep.py} & Beam divergence response as a function of screen grid extraction potential ($V_{\mathrm{s}} = 600$--$2000\,\si{\volt}$). & \texttt{divergence\_vs\_voltage\_benchmark.png} \\
\bottomrule
\end{tabular}
\caption{Summary of multi-core benchmarking scripts.}
\end{table}

% =====================================================
\section{Standalone Executable Compilation Architecture}
% =====================================================

The embedded executable builder (\texttt{Settings $\rightarrow$ Build .exe...}) runs PyInstaller asynchronously via the \texttt{PyInstallerWorker(QThread)} class. The compilation process proceeds without blocking the Qt event loop:
\begin{enumerate}[nosep]
    \item Spawns a dedicated background worker thread and builds inside a temporary directory (\texttt{\_exe\_build\_tmp/}) to keep the repository root clean.
    \item Intercepts real-time standard output and maps build milestones (analysis, dependency collection, PYZ archive creation, binary packaging) to a 0--100\% progress dialog.
    \item Automatically copies the resulting single-file executable and accompanying \texttt{config.json} into the user-selected destination directory.
    \item Cleans up intermediate build artifacts upon completion or cancellation.
\end{enumerate}

% =====================================================
\section{Summary of Execution Modes}
% =====================================================

\subsection{Interactive Graphical Run}
\begin{lstlisting}[language=bash]
python Python/main.py
\end{lstlisting}

\subsection{Headless Parameter Sweep Execution}
\begin{lstlisting}[language=bash]
python Python/run_simulation_from_config.py
\end{lstlisting}

\subsection{Parallel Multi-Core Benchmarks}
\begin{lstlisting}[language=bash]
python Python/benchmarks/benchmark_perveance.py
python Python/benchmarks/benchmark_cex.py
python Python/benchmarks/benchmark_ebs.py
\end{lstlisting}

% =====================================================
\begin{thebibliography}{9}
\bibitem{birdsall}
C.~K. Birdsall and A.~B. Langdon,
\textit{Plasma Physics via Computer Simulation},
Institute of Physics Publishing, Bristol, 1991.

\bibitem{roy}
R.~I.~S. Roy, \textit{Numerical Simulation of Ion Thruster Plume Backflow}, Ph.D. Thesis, Massachusetts Institute of Technology, Cambridge, MA, 1995.

\bibitem{matsunami}
N.~Matsunami \textit{et al.},
``Energy dependence of the yields of ion-induced sputtering of monatomic solids,'' \textit{Atomic Data and Nuclear Data Tables}, vol.~31, no.~1, pp.~1--80, 1984.
\end{thebibliography}

\end{document}
